\textbf{(CAP-COLUNI 2019 -- Q13)} Participam do campeonato brasileiro de futebol 20 equipes. Ao fim da 26ª rodada do campeonato, a classificação era a seguinte:
\begin{center}
\begin{tabular}{ll}
\multicolumn{2}{c}{Classificação}\\\hline
1º & São Paulo\\
2º & Palmeiras\\
3º & Internacional\\
4º & Flamengo\\
5º & Grêmio\\
6º & Atlético Mineiro\\
$\vdots$ & $\vdots$\\
19º & Sport\\
20º & Paraná Clube
\end{tabular}
\end{center}
\begin{center}(Disponível em: <http://globoesporte.globo.com/futebol/brasileirao-serie-a>. Acesso em: 26 set. 2018.)\end{center}

Até a 26ª rodada, o Atlético Mineiro tinha 42 pontos e o São Paulo, 35 pontos a mais que o último colocado. A pontuação que o Paraná Clube alcançou até então era $\frac{1}{3}$ da pontuação do Flamengo. Já a pontuação do Flamengo era o dobro da pontuação do Sport.

Sabendo que o Sport tem 18 pontos a menos que o Atlético Mineiro, é CORRETO afirmar que a diferença de pontos entre o primeiro e o penúltimo colocados no campeonato brasileiro é:

\begin{enumerate}
	\item 24
	\item 25
	\item 26
	\item 27
\end{enumerate}
